\vspace{-3pt}
\subsection{Training Strategy}
\label{sec3.3:training_pipeline}
\noindent{\textbf{Feed-Forward Predictive Densification.}}
To allocate a limited number of Gaussians adaptively, the densification signal must satisfy two key properties. 
First, it should correlate with potential quality gain. It should allocate more Gaussians to under-represented regions and fewer to simple regions, so that representation quality tends to improve as the number of Gaussians increases.
Second, it must be available at inference time without requiring iterative optimization. A desirable densification score must be computable using only the input images. 
To satisfy these conditions, we draw inspiration from the adaptive density control (ADC) strategy of standard 3DGS works~\cite{kerbl20233dgs, ye2024absgs}, which iteratively optimizes 3D Gaussian primitives.

\begin{figure}[t]
\centering
\begin{minipage}[t]{0.38\linewidth}
    \centering
    \vspace{5pt}
    \includegraphics[width=\linewidth]{Figures/score_loss_scheme.pdf}
    \captionsetup{type=figure}
    \vspace{-15pt}
    \caption{\textbf{Calculating $\mathcal{L}_{\mathcal{G}}^\text{score}$.} We convert the view-space rendering gradient $\mathbf{v}_g$ for each $\mathbf{g}_g\in\mathcal{G}$ into a densification score $d_g$ and learn the network to predict it.}
    \label{fig:score_loss_scheme}
\end{minipage}
\hspace{0.03\linewidth}
\begin{minipage}[t]{0.50\linewidth}
    \centering
    \vspace{5pt}
    \includegraphics[width=\linewidth]{Figures/threshold_effect.pdf}
    \captionsetup{type=figure}
    \vspace{-15pt}
    \caption{\textbf{Effect of the threshold $\tau$ across scene complexity.} For a fixed $\tau$, complex scenes yield higher densification scores and allocate more Gaussians, while simpler scenes allocate fewer.}
    \label{fig:threshold_analysis}
\end{minipage}
\vspace{-17pt}
\end{figure}

These previous works~\cite{kerbl20233dgs, ye2024absgs} periodically densify 3D Gaussians during training. In AbsGS~\cite{ye2024absgs}, for a set of predicted Gaussian primitives $\mathcal{G}=\{{\mathbf{g}_g}\}^{N_\mathcal{G}}_{g=1}$, whether to densify the Gaussian $\mathbf{g}_g$ is guided by the homodirectional target view-space positional gradient of the Gaussian, which can be obtained by backpropagating the rendering loss. 
The rendering loss $\mathcal{L}^\text{render}$ is calculated by the weighted sum of MSE and LPIPS loss between the predicted target image $\hat{\mathbf{I}}^\text{tgt}$ and the ground-truth target image $\mathbf{I}^\text{tgt}$.
With this loss, we can calculate the homodirectional view-space positional gradient $\mathbf{v}_g$:
\begin{equation}
    \mathbf{v}_g = \left( \,\sum^{m}_{j=1}\left|\frac{\partial \mathcal{L}_j^\text{render}}{\partial  \bar{\boldsymbol{\mu}}_{g,x}}\right|,\,\, \sum^{m}_{j=1}\left|\frac{\partial \mathcal{L}_j^\text{render}}{\partial \bar{\boldsymbol{\mu}}_{g,y}}\right|\,\, \right).
\end{equation}
Here $(\bar{\boldsymbol{\mu}}_{g,x}, \bar{\boldsymbol{\mu}}_{g,y})$ denotes the center of Gaussian $\mathbf{g}_g$ after projection onto the 2D image plane of the target view, $\mathcal{L}_j^\text{render}$ denotes the rendering loss computed by the $j$-th pixel of the image $\hat{\mathbf{I}}^\tet$, and $m$ is the total number of pixels that Gaussian $\mathbf{g}_g$ participates in rendering of $\hat{\mathbf{I}}^\tet$.
A large value of $\mathbf{v}_g$ indicates that the Gaussian $\mathbf{g}_g$ significantly affects the rendering loss, implying that the corresponding region is underrepresented. AbsGS~\cite{ye2024absgs} empirically shows that assigning more Gaussians based on the norm of this gradient improves the fidelity of the representation. Therefore, this gradient-based value is a suitable signal for our spatially adaptive Gaussian allocation pipeline. However, feed-forward 3DGS predicts 3D Gaussian primitives in a single forward pass without iterative scene optimization, and ground-truth images are not available at inference time. As a result, the gradient $\mathbf{v}_g$ cannot be utilized directly at inference time. 


To overcome this limitation, we instead \textit{learn to predict} $\hat{d}_g$ from the gradient $\mathbf{v}_g$. Here, $\hat{d}_g$ denotes the predicted densification score associated with the spatial region occupied by Gaussian $\mathbf{g}_g$. During the training process, using the 3D Gaussian representation $\mathcal{G}$ predicted by input context images $\{\mathbf{I}^\tec_i\}^{N_\tec}_{i=1}$, we can compute the rendering loss $\mathcal{L}^{\text{render}}(\hat{\mathbf{I}}^\tet,\mathbf{I}^\tet)$ and the homodirectional view-space positional gradient $\mathbf{v}_g$ for each Gaussian $\mathbf{g}_g\in\mathcal{G}$, as explained above. Based on the accumulated gradient $\mathbf{v}_g$, we define the supervision signal $d_g$ for the densification as a log-scaled value of the $\ell_2$ norm of the gradient: $d_g \!=\! \log\left(1+10^4\cdot \|\mathbf{v}_g\|_2\right)$. Then, the predicted densification score $\hat{d}_g$ for $\mathbf{g}_g$ is trained to match $d_g$ via the following $\ell_1$ loss:
\begin{equation}
    \mathcal{L}^{\text{score}}_\mathcal{G} = \mathbb{E}_{\mathbf{g}_g \in \mathcal{G}} \left[\|\hat{d}_g - d_g\|_1\right].
\end{equation}
The schematic illustration of this process is given in~\cref{fig:score_loss_scheme}. We use a novel view as a target view during training, allowing the model to learn densification scores that generalize better across viewpoints. Since the densification score is learned from gradient signals, it serves as an absolute and comparable criterion across different scenes, rather than a purely relative ranking within a single scene. Consequently, the threshold $\tau$ serves as a control knob for reconstruction fidelity, determining the level of detail preserved in the final representation. As shown in~\cref{fig:threshold_analysis}, under a fixed $\tau$, a complex scene generally produces higher densification scores, resulting in higher Gaussian allocation. In contrast, a simple scene yields lower scores and consequently fewer Gaussians. 

\noindent{\textbf{Training with Novel Views.}} Training the model by supervising the final 3D Gaussian representation using the input context views, rather than novel views, can lead to trivial reconstruction and potential overfitting, as the model may simply memorize the observed views instead of learning a representation that generalizes across viewpoints. To address this issue, unlike the previous feed-forward approach~\cite{jiang2025anysplat} that supervises reconstruction of input context views, we instead use a novel view as the target image during training. 
However, directly using the ground-truth camera parameters of the target view is problematic because the camera coordinate system predicted by the model may differ from the ground-truth coordinate frame. To resolve this mismatch, we align the ground-truth target camera into the predicted coordinate system before rendering the target view.
Specifically, let $\mathbf{T}^\tec_{n_1}$ and $\mathbf{T}^\tec_{n_2}$ denote the ground-truth poses of the two context views closest to the target view, and $\hat{\mathbf{T}}^\tec_{n_1}$ and $\hat{\mathbf{T}}^\tec_{n_2}$ denote their corresponding predicted poses. 
We estimate a similarity transformation matrix $\mathbf{A} \in \mathrm{Sim}(3)$
that maps poses from the ground-truth coordinate frame to the predicted one.
The context view $n_1$ closest to the target view is used as the anchor reference for rotation and translation alignment, while the scale is estimated from the relative translation between the two context views $n_1$ and $n_2$ closest to the target view.
The ground-truth target pose $\mathbf{T}^\tet$ is then transformed using the matrix $\mathbf{A}$ to obtain the aligned target pose, $\hat{\mathbf{T}}^\tet = \mathbf{A} \mathbf{T}^\tet$. For the intrinsic parameters $f$, we adjust the focal length of the target view using the ratio between the predicted and ground-truth focal lengths of the nearest context view, $\hat{f}^\tet = \hat{f}^\tec_{n_1}/ f^\tec_{n_1} \cdot f^\tet$. This simple yet carefully designed alignment procedure ensures that the aligned target view parameters $(\hat{\mathbf{T}}^\tet, \hat{\mathbf{K}}^\tet)$ are geometrically consistent with the coordinate system predicted by the model, enabling stable and reliable novel view training.

\noindent{\textbf{Overall Training Pipeline.}}
During training, we randomly sample a threshold $\tau$ and predict the corresponding Gaussian set $\mathcal{G}_\tau$, which is supervised using the rendering loss $\mathcal{L}^{\text{render}}$. In addition, we also apply the rendering loss to intermediate Gaussian representations $\mathcal{G}^l\in\mathbb{R}^{d_\mathcal{G}\times N_\tec H^l W^l}$, which are constructed using only the Gaussians predicted at each level $l$.
The densification score loss $\mathcal{L}^{\text{score}}$ is computed using the densification scores derived from the multi-level Gaussian representations $\{\mathcal{G}^l\}_{l=1}^{L-1}$. The loss for camera parameter estimation, $\mathcal{L}^{\text{camera}}$, is learned in the same way as VGGT~\cite{wang2025vggt}.
We further introduce a scene-scale regularization loss $\mathcal{L}^{\text{scene}}$, which normalizes the average distance of Gaussian centers from the origin to $1$, i.e., $\mathcal{L}^{\text{scene}} = \left|\frac{1}{|\mathcal{G}|}\sum_{\mathbf{g}_g\in\mathcal{G}} \|\boldsymbol{\mu}_g\|_2 - 1\right|$. This regularization is applied independently to the representation of each level. Without this regularization, training can become unstable. Finally, the total training objective is defined as the weighted sum of $\mathcal{L}^\text{render}$, $\mathcal{L}^\text{score}$, $\mathcal{L}^\text{camera}$, and $\mathcal{L}^\text{scene}$.
